\documentclass{article}
\usepackage{amsmath}
\usepackage{biblatex}
\addbibresource{../sources.bib}
\begin{document}

\section*{Review: Implementation of the genetic algorithm by means of CUDA technology involved in travelling salesman problem  - Anna Plichta, Tomasz Ga¸ciarz, Bartosz Baranows, Szymon Szominski 2014}

\item \textbf{Problem Context:} TSP is a well-known NP-hard problem involving finding the shortest Hamiltonian cycle visiting all cities exactly once. GA is chosen due to its effectiveness for combinatorial optimization \cite{plichta2014cuda}.

\item \textbf{GA Methodology:}
\begin{itemize}
    \item Chromosome encoding uses integer permutation representing city visiting sequences.
    \item Genetic operators include tournament and roulette wheel selection (with tournament selected for efficiency).
    \item Three crossover methods tested: one-point, two-point, and greedy crossover adapted to maintain valid TSP permutations by corrective mechanisms.
    \item Mutation swaps two random genes to maintain diversity, applied with low probability (~0.5\%).
    \item Fitness function measures route length, stored as a 2D distance matrix.
    \item Heuristics: 2-opt local optimization incorporated to improve convergence.
\end{itemize}

\item \textbf{CPU Implementation:}
\begin{itemize}
    \item Developed in C with OpenMP for multithreading.
    \item Island model dividing population into subpopulations, each evolving independently with migration of best individuals.
    \item Multiple versions evolved to include tournament selection, greedy crossover, heuristics, and island parallelization.
\end{itemize}

\item \textbf{GPU Implementation:}
\begin{itemize}
    \item Developed in CUDA C under Visual Studio 2008 with NVIDIA CUDA Toolkit 4.0.
    \item High parallelism achieved by assigning one thread per individual; requiring large population sizes (minimum ~5760 individuals for efficient GPU use).
    \item Population divided into many islands, each handled as CUDA thread blocks with communication via global memory; individual data stored in shared memory inside blocks for performance.
    \item Kernels designed for fitness evaluation, selection, crossover, mutation, and migration.
\end{itemize}

\item \textbf{Setup:}
\begin{itemize}
    \item CPU: Intel Core 2 Duo E8400 @ 3.6 GHz, dual-core.
    \item GPU: NVIDIA GeForce GTX 285 with 240 cores, 1 GB GDDR3 memory at 702 MHz.
    \item Time measurements done with QueryPerformanceCounter on CPU and cudaEventRecord on GPU.
    \item Chromosome lengths (number of cities) tested up to 100.
\end{itemize}

\item \textbf{Performance Results:}
\begin{itemize}
    \item GPU implementation shows near-constant execution time even when increasing population size, unlike linearly increasing times on CPU.
    \item Speedup up to 23x relative to single CPU core, and 13x relative to dual-core CPU.
    \item Demonstrates that large-scale parallelism on GPU effectively leverages CUDA capabilities for GAs solving TSP.
    \item Performance gains require sufficiently large population sizes (thousands of individuals).
\end{itemize}

\item \textbf{Significance:} The paper validates the substantial acceleration achieved by adapting GA to GPU via CUDA, emphasizing design considerations for GPU computing such as massive parallelism, memory hierarchy, and appropriate problem/algorithm granularity for efficiency \cite{plichta2014cuda}.

\bibliographystyle{plain}
\bibliography{sources}

\end{document}